\section{Setup}
\label{sec:setup}

\subsection{Triplet Schema}
\label{sec:triplets}

We start from WikiText-103 \citep{merity2017pointer}, filter to $16{,}038$ usable passages, and embed them with \texttt{all-MiniLM-L6-v2} \citep{wang2020minilmv2} into $384$-dim vectors. HDBSCAN \citep{campello2013density} with \texttt{min\_cluster\_size}{=}$200$, \texttt{min\_samples}{=}$5$, Euclidean metric, and EOM cluster selection produces $10$ semantic clusters (covering recurring WikiText topics: Federer, Harry Potter, tennis matches, etc.) plus a noise pool of $9{,}198$ ($57.4\%$).

For each cluster we sample $10$ disjoint \emph{triplets}, giving $100$ triplets in total. Each triplet contains three same-cluster but mutually disjoint splits of $50$ texts:
\begin{center}
\small
\begin{tabular}{@{}lll@{}}
    \toprule
    Split & 50 texts & Downstream role \\
    \midrule
    \texttt{train}      & forget set $\Df$              & fed to GradAscent \\
    \texttt{validation} & retain neighbours             & retain loss + $\Ltwo$ probe \\
    \texttt{test}       & probe samples                 & $\Ltwo / \Lthree$ evaluation \\
    \bottomrule
\end{tabular}
\end{center}
The within-cluster but mutually disjoint design lets us read off three semantic distances by simply re-indexing the same evaluation matrix. Different triplets may overlap across clusters (sampling is independent), but within a triplet the three splits never share a text.

\subsection{Unlearning Procedure}
\label{sec:unlearn}

We unlearn Llama-3.1-8B-Instruct \citep{dubey2024llama3} with the GradAscent objective \citep{jang2023knowledge}: for each $\Df$ we maximise the negative log-likelihood on $\Df$ while keeping a small retain-loss term on the same-cluster validation split. Hyperparameters follow TOFU-aligned defaults from OpenUnlearning \citep{maini2024tofu, dorna2025openunlearning}: learning rate $1\!\times\!10^{-5}$ with linear decay, paged AdamW (32-bit), per-device batch size $8$ with gradient accumulation $8$ (effective batch $64$), $5$ epochs (resulting in $5$ optimiser steps because $|\Df|{=}50<64$), warmup of $1$ step, weight decay $0.01$, BF16. All runs use a single H100 80\,GB. Each unlearn run takes $\approx 75$\,s wall-clock; the full $n{=}100$ sweep finishes in $1$\,h\,$54$\,min and produces $\approx 1.5$\,TB of checkpoints ($\approx 15$\,GB each). Train loss converges in $[-0.84, -0.62]$ (negative is expected for GradAscent).

\subsection{Cross-PPL Matrix and Three Layers}
\label{sec:matrix}

Given $100$ unlearned checkpoints and $100$ evaluation triplets, we evaluate every checkpoint on every triplet's three splits, yielding $100\times100\times3\times50 = 1{,}500{,}000$ per-sample perplexities (Fig.~\ref{fig:cross-matrix}; numerically: $5{,}000$ $\Lone$ samples on the diagonal train splits, $10{,}000$ $\Ltwo$ samples on diagonal val/test, and $495{,}000$ $\Lthree$ samples on the off-diagonal test). We pair each measurement with the corresponding base-model PPL and define
\begin{equation}
    r_{m,e,s,i} \;=\; \frac{\PPL_{\text{unlearned},m}(x_i^{e,s})}{\PPL_{\text{base}}(x_i^{e,s})},
\end{equation}
where $m$ indexes the unlearned checkpoint, $e$ the evaluation triplet, $s\in\{\text{train},\text{val},\text{test}\}$ the split, and $i\in[50]$ the sample. The three layers correspond to:
\begin{align*}
    \Lone&: m{=}e,\, s{=}\text{train} \quad\text{(forget)}\\
    \Ltwo&: m{=}e,\, s\in\{\text{val},\text{test}\} \quad\text{(locality)}\\
    \Lthree&: m{\neq}e,\, s{=}\text{test} \quad\text{(spillover)}
\end{align*}
We summarise each layer by the geometric mean of $r$, equivalently $\exp(\overline{\log r})$, and the fraction of samples with $r{>}1.1$.

The full sweep takes $\approx 6$\,h\,$11$\,min on H100 ($\approx 2.2$\,s per (model, eval) pair). We additionally cache the per-sample long-format table (\texttt{ppl\_long.parquet}, $1.5$M rows) for downstream auditing and four invariant sanity checks (Sec.~\ref{sec:three-layer}).

\subsection{Audit Features}
\label{sec:features}

For each forget set $\Df$ of $50$ texts, we compute a $12$-dim intrinsic geometric descriptor from the same MiniLM embeddings used for clustering, with no access to model weights, gradients, or any unlearned checkpoint:
\begin{itemize}
    \setlength\itemsep{0.15em}
    \item \textbf{Spread (3)}: \texttt{emb\_variance\_mean}, \texttt{emb\_variance\_max}, \texttt{pairwise\_eucl\_mean}.
    \item \textbf{Similarity (3)}: \texttt{pairwise\_sim\_mean / std / q90} (cosine).
    \item \textbf{Location/scale (3)}: \texttt{centroid\_norm}, \texttt{emb\_norm\_mean / std}.
    \item \textbf{Shape (3)}: \texttt{effective\_rank}, \texttt{isotropy}, \texttt{spread\_over\_centroid}.
\end{itemize}
All $12$ features are scalars derived from the $50\times 384$ forget-set embedding matrix; computing them takes milliseconds per $\Df$.
